\subsection{工学微積分}
\term{微積分}{Calculus}は、連続的な変化を扱う数学である。工学微積分では、極限、連続性、微分、積分、超越関数、ベクトル解析を、工学上の問題に適用する。

\par\smallskip\noindent\refstepcounter{table}\label{tab:ch17-21}表 \thetable: 工学微積分における概念・説明と用途\par\smallskip
表 \thetable は、工学微積分における概念・説明と用途を比較・確認しやすい形で整理したものである。\par\smallskip
\begin{longtable}{>{\raggedright\arraybackslash}p{0.26\linewidth} >{\raggedright\arraybackslash}p{0.68\linewidth}}
\toprule
概念 & 説明と用途 \\
\midrule
\endfirsthead
\toprule
概念 & 説明と用途 \\
\midrule
\endhead
極限 & 入力がある値に近づくとき、関数値がどこへ近づくかを見る。 \\
連続性 & 小さな入力変化が小さな出力変化に対応する性質である。 \\
微分 & 変化率を扱う。最適化、速度、勾配、感度分析に使う。 \\
積分 & 面積、累積量、総量を扱う。負荷の累積、エネルギー、確率密度に使う。 \\
超越関数 & 指数関数、対数関数、三角関数など、代数方程式だけでは表せない関数である。 \\
ベクトル解析 & 三次元空間のベクトル場の微分・積分を扱う。物理、制御、画像処理に関係する。 \\
\bottomrule
\end{longtable}
\begin{notebox}{ソフトウェアでの見方}
微積分は、機械学習の最適化、物理シミュレーション、制御、画像処理、信号処理、性能の傾向分析で使われる。初学者は、公式暗記よりも「変化率」「累積量」「近似」という意味を理解することが重要である。
\end{notebox}
\par\smallskip
\begin{center}
\centering
\IfFileExists{assets/generated_figures/ch17/fig-ch17-16.png}{\includegraphics[width=0.96\linewidth]{assets/generated_figures/ch17/fig-ch17-16.png}}{\fbox{\parbox[c][48mm][c]{0.9\linewidth}{\centering 画像生成待ち\\微分と積分の用途図}}}
\par\smallskip
\refstepcounter{figure}\label{fig:ch17-16}図 \thefigure: 微分と積分の用途図
\end{center}
図 \thefigure は、微分と積分の用途図を視覚的に整理したものである。
\par\smallskip
